\chapter{Analyse et specification des besoins}
\label{chap:analyse}

\section{Identification des acteurs}

L'analyse du perimetre fonctionnel (\cref{sec:perimetre}) fait apparaitre trois
acteurs interagissant avec la plateforme :

\begin{itemize}
  \item \textbf{Analyste securite / RSSI} : consulte le tableau de bord, explore le
  detail du score d'un actif, interroge l'assistant conversationnel. C'est
  l'utilisateur principal du systeme au quotidien.
  \item \textbf{Administrateur fonctionnel} : en plus des actions de l'analyste,
  importe les nouveaux exports Qualys/CMDB et configure les parametres du modele de
  risque (poids, facteurs, seuils SLA). Dans le modele UML, cet acteur
  \emph{generalise} l'analyste : tout ce que l'analyste peut faire, l'administrateur
  le peut aussi.
  \item \textbf{Moteur LLM local (Ollama)} : acteur secondaire, externe au systeme,
  sollicite par le cas d'utilisation « Interroger l'assistant RAG » pour generer une
  reponse en langage naturel a partir des donnees recuperees par les outils exposes
  a l'agent.
\end{itemize}

\section{Besoins fonctionnels}

\begin{longtable}{@{}p{1.3cm}p{11.5cm}@{}}
\toprule
\textbf{Code} & \textbf{Besoin fonctionnel} \\
\midrule
\endhead
BF-01 & Importer un fichier d'export Qualys (XML) et en extraire la liste des hotes
et de leurs vulnerabilites. \\
BF-02 & Importer un fichier d'export CMDB (XML) et en extraire la liste des elements
de configuration (CI). \\
BF-03 & Rejeter tout fichier deja traite (detection par empreinte cryptographique) et
tout fichier dont le type MIME reel ne correspond pas au type declare. \\
BF-04 & Recalculer automatiquement, des qu'un import aboutit ou qu'une configuration
change, l'ACS, le score QDS pondere, l'ARS, la classe de risque et le statut SLA de
chaque actif reconcilie. \\
BF-05 & Permettre a l'administrateur de configurer les poids de severite
($w_c, w_h, w_m, w_l$), les poids par role technique et par environnement, les
facteurs de modulation ($I_{External}$, $M_{SOC}$, $M_{AV}$, $M_{Patch}$) et les
seuils/delais SLA par classe de risque. \\
BF-06 & Marquer un actif en exception (« risque accepte ») et refleter cet etat dans
le statut SLA affiche, sans alterer l'ARS calcule. \\
BF-07 & Afficher un tableau de bord synthetisant la repartition du parc par classe de
risque, par equipe/proprietaire, les flux Environnement/Role/Classe, l'historique de
l'ARS moyen et la conformite SLA. \\
BF-08 & Mettre a jour le tableau de bord en temps reel (sans rechargement de page)
des qu'un recalcul de risque aboutit. \\
BF-09 & Afficher, pour un actif donne, le detail explicatif du calcul de son ARS
(contribution de chaque terme de la formule). \\
BF-10 & Permettre a l'analyste d'interroger en langage naturel les donnees ingerees
et calculees via un assistant conversationnel. \\
\bottomrule
\end{longtable}

\section{Besoins non fonctionnels}

\begin{longtable}{@{}p{2.6cm}p{10.2cm}@{}}
\toprule
\textbf{Categorie} & \textbf{Exigence} \\
\midrule
\endhead
Performance & Le calcul de l'ARS pour l'ensemble du parc doit s'effectuer en
parallele (un couple ACS/QDS pondere par actif) pour rester compatible avec des
imports comportant plusieurs milliers d'hotes. \\
Reactivite & Le tableau de bord doit refleter un nouveau recalcul sans action de
l'utilisateur (push serveur $\to$ client). \\
Coherence transactionnelle & Un evenement declenchant un recalcul ne doit etre traite
qu'apres la validation (\emph{commit}) de la transaction qui l'a emis, afin d'eviter
de lire un etat de base de donnees partiel. \\
Idempotence des imports & Un meme fichier ne doit jamais etre traite deux fois ; un
nouvel import d'un meme type de source remplace integralement l'etat precedent. \\
Robustesse des entrees & Le type reel d'un fichier uploade doit etre verifie
independamment de son extension declaree (protection contre l'upload de fichiers non
maitrise). \\
Configurabilite & Aucun parametre du modele de risque (poids, facteurs, seuils) ne
doit etre code en dur : tous doivent rester modifiables sans redeploiement. \\
Portabilite & L'application doit pouvoir etre deployee de maniere reproductible via
conteneurisation (Docker Compose), y compris le moteur d'inference du LLM local. \\
\bottomrule
\end{longtable}

\section{Diagramme de cas d'utilisation}

Le \cref{fig:uc-diagram} synthetise les besoins fonctionnels precedents sous forme de
diagramme de cas d'utilisation UML. Les relations \texttt{<<include>>} traduisent le
fait que trois evenements distincts (import Qualys, import CMDB, soumission d'une
configuration) declenchent systematiquement le meme sous-processus : le recalcul du
risque.

\begin{figure}[H]
\centering
\begin{tikzpicture}[
  actor/.style={draw,minimum width=1.6cm,align=center,font=\small,inner sep=6pt},
  usecase/.style={draw,ellipse,minimum width=3.6cm,minimum height=1cm,align=center,font=\scriptsize,fill=dxcblue!6},
  include/.style={draw,dashed,-{Latex[length=2mm]}},
  assoc/.style={draw,-},
  gen/.style={draw,-{Triangle[length=2.6mm,width=2.6mm,open]}}
]

% Acteurs
\node[actor] (analyste) at (-0.2,-2.3) {
  \begin{tikzpicture}[scale=0.35]
    \draw (0,0) circle (0.4);
    \draw (0,-0.4) -- (0,-1.6);
    \draw (-0.7,-0.7) -- (0.7,-0.7);
    \draw (0,-1.6) -- (-0.6,-2.5);
    \draw (0,-1.6) -- (0.6,-2.5);
  \end{tikzpicture}\\ Analyste securite};

\node[actor] (admin) at (-0.2,-7.4) {
  \begin{tikzpicture}[scale=0.35]
    \draw (0,0) circle (0.4);
    \draw (0,-0.4) -- (0,-1.6);
    \draw (-0.7,-0.7) -- (0.7,-0.7);
    \draw (0,-1.6) -- (-0.6,-2.5);
    \draw (0,-1.6) -- (0.6,-2.5);
  \end{tikzpicture}\\ Administrateur\\ fonctionnel};

\node[actor] (llm) at (13.4,-6.6) {
  \begin{tikzpicture}[scale=0.35]
    \draw (0,0) circle (0.4);
    \draw (0,-0.4) -- (0,-1.6);
    \draw (-0.7,-0.7) -- (0.7,-0.7);
    \draw (0,-1.6) -- (-0.6,-2.5);
    \draw (0,-1.6) -- (0.6,-2.5);
  \end{tikzpicture}\\ Moteur LLM\\ (Ollama)};

% Frontiere du systeme
\node[draw,minimum width=9.6cm,minimum height=9.6cm,label={[font=\bfseries\small]above:Plateforme}] (sys) at (5.7,-4.2) {};

% Cas d'utilisation
\node[usecase] (uc1) at (2.3,-0.5)  {Importer un rapport\\ Qualys};
\node[usecase] (uc2) at (2.3,-2.0)  {Importer un export\\ CMDB};
\node[usecase] (uc3) at (2.3,-3.5)  {Configurer les poids\\ et facteurs du modele};
\node[usecase] (uc4) at (2.3,-5.0)  {Configurer les seuils\\ et delais SLA};
\node[usecase] (uc5) at (2.3,-6.5)  {Marquer un actif\\ en exception};
\node[usecase] (uc6) at (8.9,-1.0)  {Recalculer le risque\\ (ARS, classe, SLA)};
\node[usecase] (uc7) at (8.9,-2.6)  {Consulter le tableau\\ de bord temps reel};
\node[usecase] (uc8) at (8.9,-4.1)  {Consulter le detail\\ explicatif d'un actif};
\node[usecase] (uc9) at (8.9,-5.6)  {Interroger l'assistant\\ conversationnel (RAG)};
\node[usecase] (uc10) at (8.9,-7.1) {Suivre l'historique\\ de l'ARS moyen};

% Associations acteurs - cas d'utilisation
\draw[assoc] (analyste) -- (uc7);
\draw[assoc] (analyste) -- (uc8);
\draw[assoc] (analyste) -- (uc9);
\draw[assoc] (analyste) -- (uc10);

\draw[assoc] (admin) -- (uc1);
\draw[assoc] (admin) -- (uc2);
\draw[assoc] (admin) -- (uc3);
\draw[assoc] (admin) -- (uc4);
\draw[assoc] (admin) -- (uc5);

\draw[gen] (admin) -- (analyste);

\draw[assoc] (uc9) -- (llm);

% Include
\draw[include] (uc1) -- (uc6) node[midway,above,sloped,font=\tiny]{<<include>>};
\draw[include] (uc2) -- (uc6) node[midway,above,sloped,font=\tiny]{<<include>>};
\draw[include] (uc3) -- (uc6) node[midway,below,sloped,font=\tiny]{<<include>>};
\draw[include] (uc4) -- (uc6) node[midway,below,sloped,font=\tiny]{<<include>>};

\end{tikzpicture}
\caption{Diagramme de cas d'utilisation de la plateforme}
\label{fig:uc-diagram}
\end{figure}

\section{Description textuelle des cas d'utilisation principaux}

\subsection{UC -- Importer un rapport Qualys / un export CMDB}

\begin{longtable}{@{}p{3.2cm}p{9.6cm}@{}}
\toprule
\textbf{Acteur principal} & Administrateur fonctionnel \\
\textbf{Pre-condition} & Le fichier est un document XML dont le contenu correspond au
format attendu (rapport Qualys ou export CMDB). \\
\textbf{Scenario nominal} &
1) L'administrateur televerse le fichier via l'interface d'upload. \\
& 2) Le systeme detecte le type MIME reel du fichier et calcule son empreinte
SHA. \\
& 3) Si l'empreinte correspond a un fichier deja traite, l'import est rejete. \\
& 4) Sinon, le fichier est stocke sur disque et marque comme en attente de
traitement. \\
& 5) L'administrateur declenche le lancement du traitement par lots
(\emph{batch}) correspondant au type de fichier. \\
& 6) Le systeme vide les donnees precedentes de ce type de source, puis lit,
transforme et persiste les nouvelles donnees par lots. \\
& 7) A l'issue du traitement, un evenement d'ingestion terminee est publie,
declenchant le recalcul du risque (cas d'utilisation inclus). \\
\textbf{Post-condition} & Les donnees de la source concernee refletent
integralement le nouvel export ; le risque de tous les actifs impactes est
recalcule. \\
\textbf{Exceptions} & Fichier deja traite (empreinte identique) ; type MIME
incoherent avec le type declare ; traitement deja en cours (verrou applicatif). \\
\bottomrule
\end{longtable}

\subsection{UC -- Consulter le tableau de bord temps reel}

\begin{longtable}{@{}p{3.2cm}p{9.6cm}@{}}
\toprule
\textbf{Acteur principal} & Analyste securite \\
\textbf{Scenario nominal} &
1) L'analyste ouvre la page du tableau de bord. \\
& 2) Le systeme calcule et renvoie l'ensemble des agregats (repartition par
classe, par equipe, flux Sankey, historique, points SLA). \\
& 3) Le navigateur ouvre une connexion evenementielle (SSE) vers le serveur. \\
& 4) A chaque recalcul de risque abouti, le serveur pousse les nouveaux agregats
a tous les clients connectes, qui redessinent les graphiques sans rechargement
de page. \\
\textbf{Post-condition} & L'analyste dispose d'une vue a jour de la posture de
risque du parc. \\
\bottomrule
\end{longtable}

\subsection{UC -- Interroger l'assistant conversationnel (RAG)}

\begin{longtable}{@{}p{3.2cm}p{9.6cm}@{}}
\toprule
\textbf{Acteur principal} & Analyste securite \\
\textbf{Acteur secondaire} & Moteur LLM local (Ollama) \\
\textbf{Scenario nominal} &
1) L'analyste pose une question en langage naturel (ex. : « quels sont les
actifs en classe Black non couverts par le SOC ? »). \\
& 2) L'agent conversationnel determine si la question necessite d'interroger la
base de donnees via un outil expose (\emph{tool calling}). \\
& 3) Si oui, l'outil est invoque, son resultat est injecte dans le contexte du
modele. \\
& 4) Le modele local genere une reponse en flux (\emph{streaming}) restituee
progressivement a l'utilisateur. \\
\textbf{Post-condition} & L'analyste obtient une reponse fondee sur les donnees
reelles du systeme plutot que sur la seule connaissance generale du modele. \\
\bottomrule
\end{longtable}
